\documentclass[11pt]{article}

\usepackage[margin=1in]{geometry}
\usepackage[T2A]{fontenc}
\usepackage[utf8]{inputenc}
\usepackage[english,russian]{babel}
\usepackage{amsmath}
\usepackage{amssymb}
\usepackage{booktabs}

\title{Модулярная параметризация отношений масс лептонов и адронов в представлении по основанию $e$}
\author{И. Михайлов}
\date{Май 2026}

\begin{document}
\maketitle

\begin{abstract}
В работе представлена компактная феноменологическая параметризация выбранных
отношений масс лептонов и адронов через константы $e$, $\pi$, обратную
постоянную тонкой структуры $\alpha^{-1}$, топологический индекс $N=5$,
референсную размерность бозонной струны $D=26$ и число Фибоначчи
$F_{26}=121393$. Конструкция не выводится из лагранжиана и не предлагается как
замена Стандартной модели. Ее более узкая цель состоит в фиксации
низкосложностного набора алгебраических формул, их численных невязок и точек,
в которых анзац либо не работает, либо требует нестандартной интерпретации.
\end{abstract}

\section{Дискретный базис}

Рабочий базис задается как
\begin{equation}
    P=\{2,3,5,13,137\}, \qquad N=5,\qquad D=26,\qquad F_{26}=121393.
\end{equation}
Целое число $137$ используется как дискретный якорь. В исследовательской
записи по основанию $e$, использованной в этом проекте,
\begin{equation}
    137_e = 21100.0210102\ldots_e ,
\end{equation}
а первый фрагмент дробного цифрового паттерна имеет вид
\begin{equation}
    d_{-1},d_{-2},d_{-3},d_{-4},d_{-5},d_{-6},d_{-7}
    =
    0,2,1,0,1,0,2.
\end{equation}
Ранее использованный цифровой момент
\begin{equation}
    I_5=\sum_{k=1}^{5}k^2d_{-k}=42
\end{equation}
имеет эквивалентную топологическую форму
\begin{equation}
    I_5=2(D-N)=2(26-5)=42.
\end{equation}
Основной оператор сжатия записывается как
\begin{equation}
    q = \frac{N(N-2)}{e^4\pi^3}
      = \frac{15}{e^4\pi^3}
      =0.008860611874290873.
\end{equation}
Также используется
\begin{equation}
    s=\frac{\alpha}{2\pi}, \qquad
    \delta\phi=\cos(1/N)-\cos(2/N)=0.05900558383835652.
\end{equation}

\section{Сектор заряженных лептонов}

Электрон используется как нормировочный якорь. Проектная нормировка задается
выражением
\begin{equation}
    C_e=\frac{1}{1-s\left(N+\frac{2}{eN}\right)},
\end{equation}
а масштаб электрона записывается как
\begin{equation}
    m_e = 10e^{11}\cos^2(2/N)C_e .
\end{equation}
В данной заметке это уравнение является нормировочной конвенцией.

Отношение массы мюона к массе электрона задается как
\begin{equation}
    \frac{m_\mu}{m_e}
    =
    \alpha^{-1}\left(\frac{3}{2}+q\right).
\end{equation}
Отношение массы тау-лептона к массе электрона записывается как
\begin{equation}
    \frac{m_\tau}{m_e}
    =
    \alpha^{-1}\left[N^2+qf_\tau\right],
\end{equation}
где
\begin{equation}
    f_\tau = 2(D-N)+2e^{-2}+2e^{-7}.
\end{equation}

\section{Адронные и барионные benchmark-формулы}

Для нейтрального пиона базовое выражение имеет вид
\begin{equation}
    \frac{m_{\pi^0}}{m_e}
    =
    \alpha^{-1}
    \left[
        2-\delta\phi
        -
        q\left(1+\frac{1}{N}+\frac{1}{\pi}\right)
    \right].
\end{equation}
Аудит Phase 9 показал, что буквальная замена $q\rightarrow117/137$ не проходит
проверку, поскольку разрушает масштаб выражения. Сохраненная компактная
модуляция использует инвариант
\begin{equation}
    117=DN-13=9\cdot13
\end{equation}
только как малую поправку:
\begin{equation}
    q_{117}=q\left(1-\frac{2^2}{117\cdot137}\right).
\end{equation}
Итоговая benchmark-формула для пиона поэтому имеет вид
\begin{equation}
    \frac{m_{\pi^0}}{m_e}
    =
    \alpha^{-1}
    \left[
        2-\delta\phi
        -
        q_{117}\left(1+\frac{1}{N}+\frac{1}{\pi}\right)
    \right].
\end{equation}

Для отношения массы протона к массе электрона низкосложностное выражение Phase
9 записывается как
\begin{equation}
    K_p=\frac{M_p}{m_e}
    =
    \frac{2F_{26}}{DN}\cos(1/N)
    \left[
        1+es+\frac{6}{137^2(D+N+2^2)}
    \right].
\end{equation}
Целое $6$ является зазором функции Эйлера:
\begin{equation}
    \phi(137)-DN=136-130=6.
\end{equation}
Эта формула должна читаться как компактный эмпирический кандидат, а не как
вывод из КХД.

\section{Стерильные benchmark-предсказания}

Экстраполяция заряженного тяжелого лептона задается выражением
\begin{equation}
    \frac{m_{\tau'}}{m_e}
    =
    \alpha^{-1}\left[D^2+q(DN)\right],
\end{equation}
что дает
\begin{equation}
    m_{\tau'}=47.417730830364825~\mathrm{GeV}.
\end{equation}
Это значение выше проектного референсного порога $45~\mathrm{GeV}$, но ниже
стандартного ограничения для последовательного тяжелого заряженного лептона,
порядка $100.8~\mathrm{GeV}$. Поэтому оно не может интерпретироваться как
обычный четвертый заряженный лептон Стандартной модели без дополнительных
допущений.

Стерильный нейтральный benchmark задается выражением
\begin{equation}
    m_{\nu_{\tau'}}
    =
    m_e
    \left(\frac{F_{26}}{\alpha^{-1}}\right)
    \left[
        1+\frac{DN}{\pi}\delta\phi(1-s)
    \right],
\end{equation}
что дает
\begin{equation}
    m_{\nu_{\tau'}}=1.5566463427697075~\mathrm{GeV}.
\end{equation}
Такая масса кинематически относится к холодному сектору, однако реликтовая
плотность, время жизни, история рождения и видимые каналы распада не
определяются одной формулой массы. Benchmark условно жизнеспособен только для
полностью стерильного или достаточно сильно развязанного состояния. Простой
аудит Fermi-LAT/Galactic Center не поддержал прямую интерпретацию через
гамма-избыток: фазово-модулированные оценки аннигиляции существенно ниже
номинального масштаба $\sim2\times10^{-26}~\mathrm{cm^3\,s^{-1}}$.

\section{Численная сводка}

\begin{table}[h]
\centering
\caption{Benchmark-отношения масс, зафиксированные моделью. Невязки указаны
относительно референсных значений, использованных в проектном аудите.}
\begin{tabular}{llll}
\toprule
Observable & Formula family & Model value & Residual \\
\midrule
$m_\mu/m_e$ & $\alpha^{-1}(3/2+q)$ & $206.76822$ & $-0.29~\mathrm{ppm}$ \\
$m_\tau/m_e$ & $\alpha^{-1}(N^2+qf_\tau)$ & $3477.22820$ & $-0.022~\mathrm{ppm}$ \\
$m_{\pi^0}$ & $q_{117}$ pion branch & $134.976797~\mathrm{MeV}$ & $-0.020995~\mathrm{ppm}$ \\
$M_p/m_e$ & $\phi(137)$ gap branch & $1836.152490$ & $-0.098658~\mathrm{ppm}$ \\
$m_{\tau'}$ & $\alpha^{-1}[D^2+qDN]$ & $47.417731~\mathrm{GeV}$ & benchmark \\
$m_{\nu_{\tau'}}$ & sterile branch & $1.556646~\mathrm{GeV}$ & benchmark \\
\bottomrule
\end{tabular}
\end{table}

\section{Заключение}

Собранные здесь компактные соотношения следует рассматривать как модулярную
параметризацию выбранных масштабов частиц. Формулы численно точны и
воспроизводимы в сопроводительном коде, но остаются феноменологическими до тех
пор, пока не будут связаны с динамической теорией поля, ускорительными
ограничениями и стандартными адронными расчетами. Поэтому полезный результат
упражнения состоит не в утверждении новой физики, а в формировании
фальсифицируемой матрицы низкосложностных численных benchmark-формул.

\end{document}
